
\subsection{Replace a provisional assessment without charging twice}

Later activity changes the question that the tax calculation must answer.
The continued book includes the return adjustment of $-24$, while its earlier
assessment covers only the preceding interval. A transition to the next
financial opening cannot make that unassessed activity disappear. The next
finite method therefore supplies an explicitly selected cumulative replacement:
a later calculation over the same original opening supersedes the earlier
provisional assessment. This is a hypothetical tax model. Independent tax
periods require a different method.

The replacement reads the actual continued 26-account book. It preserves the
cash already paid, retains the separate payable and includes every pretax
account and relevant lease right in the new declarations. Each native return
adjustment must appear exactly once with its event identity and amount.
Explicit additional assumptions must reconcile to the declared account totals,
and right-level tax bases must agree with their aggregate balances. Missing
declarations remain unresolved; incompatible scope, amounts or dates are refused.

The loss calculation returns to the original opening inventory. Using the
provisional closing ledger as a new opening could count a loss twice or make
it available inside the cumulative interval that generated it. The replacement
instead applies the declared usage, cap, eligibility and expiry rule once to
the original lots. The earlier result remains attached as superseded evidence.
For example, a provisional generated loss of 19 replaced by a cumulative loss
of 7 leaves 7, not 26. An original opening loss of 19, provisionally used by 1
and later used by 19 cumulatively, has total usage 19 rather than 20.

\begin{samepage}
Let $C_0$ be current-tax expense already recorded, $C_1$ the replacement total
after explicit loss usage, and $S$ actual tax payments. The additional expense
and resulting unpaid liability are
\[
 \Delta C=C_1-C_0,\qquad L_1=C_1-S.
\]
\end{samepage}
The method debits current-tax expense and credits tax payable by $\Delta C$.
A negative amount reverses those balances; it does not change cash. The
selected profile requires $L_1\geq0$. If earlier payments exceed the new total,
it preserves the calculation but refuses the financial proposal because a
tax receivable or refund needs a separately specified method.

In the paid mixed example, expense and payments are both initially $1/2$.
A declared replacement total of $21/4$ therefore adds only $19/4$ to expense
and payable. Cash stays at $721/2$. Pretax profit of 5 less total current-tax
expense of $21/4$ gives a recorded loss of $1/4$; assets of $1401/2$ less
liabilities of $435/4$ equal equity of $2367/4$. The operating cash result
remains $-59/2$. These amounts follow the explicit cumulative assumptions.

If no tax has been paid, replacing expense of $1/2$ with a total of $1/4$
reverses expense and payable by $1/4$ without moving cash. If $1/2$ was already
paid, the same replacement would leave a receivable of $1/4$ and is refused by
this profile. Both outcomes retain the complete continued financial history
and the replacement calculation for inspection.

The new coverage endpoint is later than the provisional endpoint and no later
than the original financial period boundary. The recorded events must fit
within it. Validation recomputes the continued book, the cumulative calculation,
the replacement losses and the incremental journal; applying a proposal compares
the whole result. Deferred-tax balances remain unchanged. Further assessment
chains, deferred-tax recognition, financial-period carry, source admission and
annual completeness still require their own methods and evidence.
